% =====================================================================
\chapter{Sistemas de numeración y coma flotante}
\label{tema:6}
% =====================================================================

Un ordenador solo almacena bits. Este tema estudia cómo se representan con ellos los enteros sin signo, los enteros con signo y los números reales, y qué se pierde por el camino. La representación posicional se fundamenta en la división euclídea de enteros; la parte de máquina sigue el capítulo~2 «Representing and Manipulating Information» de Bryant y O'Hallaron~\cite{csapp} y el artículo de Goldberg~\cite{goldberg}; los formatos de coma flotante son los de la norma IEEE~754~\cite{ieee754}.

\section{Representación posicional}

\begin{teorema}[Representación en base $b$]\label{teo:base}
Sea $b\ge2$ un entero. Todo $n\in\N$ se escribe de manera única como
\[
n=\sum_{k=0}^{m}d_kb^k=(d_md_{m-1}\dots d_1d_0)_b,\qquad d_k\in\{0,1,\dots,b-1\},\ d_m\neq0.
\]
Los $d_k$ son las \emph{cifras} de $n$ en base $b$.
\end{teorema}

\begin{proof}
Se usa la división euclídea de enteros: para $n\in\N_0$ existen únicos $q,r$ con $n=bq+r$ y $0\le r<b$.

\emph{Existencia}, por inducción fuerte sobre $n$. Si $n<b$, basta tomar $m=0$ y $d_0=n$. Si $n\ge b$, se escribe $n=bq+r$ con $0\le r<b$; entonces $1\le q<n$ y, por la hipótesis de inducción, $q=\sum_{k=0}^{m'}e_kb^k$. Luego $n=r+\sum_{k=0}^{m'}e_kb^{k+1}$, que es una representación de $n$ con $d_0=r$ y $d_{k+1}=e_k$.

\emph{Unicidad}, también por inducción fuerte. Si $n=\sum_kd_kb^k=\sum_kd'_kb^k$, entonces $n=b\bigl(\sum_{k\ge1}d_kb^{k-1}\bigr)+d_0$ con $0\le d_0<b$, y lo mismo con las $d'_k$. Por la unicidad del resto, $d_0=d'_0$, y por la del cociente, $\sum_{k\ge1}d_kb^{k-1}=\sum_{k\ge1}d'_kb^{k-1}$, un número menor que $n$ al que se aplica la hipótesis de inducción (si es~0, no quedan más cifras no nulas).
\end{proof}

La demostración es un algoritmo: las cifras de $n$ en base $b$ son los restos de dividir sucesivamente entre $b$, leídos de abajo arriba. En sentido contrario, el valor de $(d_m\dots d_0)_b$ se obtiene con el \emph{esquema de Horner}, $(\cdots((d_mb+d_{m-1})b+d_{m-2})\cdots)b+d_0$, que es la regla de Ruffini del \cref{tema:3} evaluando en $x=b$.

\begin{ejemplo}\label{ej:conversion}
$45=2\cdot22+1$, $22=2\cdot11+0$, $11=2\cdot5+1$, $5=2\cdot2+1$, $2=2\cdot1+0$, $1=2\cdot0+1$. Leyendo los restos de abajo arriba, $45=(101101)_2$. Comprobación por Horner: $((((1\cdot2+0)2+1)2+1)2+0)2+1=45$.
\end{ejemplo}

\begin{proposicion}[Número de cifras]\label{prop:cifras}
Un entero $n\ge1$ tiene exactamente $\lfloor\log_bn\rfloor+1$ cifras en base $b$. Con $m$ cifras en base $b$ se representan exactamente los $b^m$ enteros $0,1,\dots,b^m-1$.
\end{proposicion}

\begin{proof}
Si $n$ tiene $m+1$ cifras con $d_m\neq0$, entonces $b^m\le n\le\sum_{k=0}^{m}(b-1)b^k=b^{m+1}-1$ por la suma geométrica (\cref{teo:geometrica}). Tomando $\log_b$, $m\le\log_bn<m+1$, es decir, $m=\lfloor\log_bn\rfloor$. La segunda afirmación se sigue del \cref{teo:base} y de que $(b-1\ \dots\ b-1)_b=b^m-1$.
\end{proof}

Las cifras de la parte fraccionaria de un número $x\in[0,1)$ se obtienen multiplicando sucesivamente por $b$ y tomando la parte entera: $x=\sum_{k\ge1}d_kb^{-k}=(0{,}d_1d_2\dots)_b$.

\begin{proposicion}\label{prop:fraccion-binaria}
Sea $p/q$ una fracción irreducible con $0<p<q$. Su desarrollo en base~2 es finito si y solo si $q$ es una potencia de~2.
\end{proposicion}

\begin{proof}
Si el desarrollo tiene $k$ cifras, $p/q=N/2^k$ con $N$ entero, es decir, $p\,2^k=qN$. Entonces $q$ divide a $p\,2^k$ y, como no tiene factores comunes con $p$, divide a $2^k$, luego es potencia de~2. Recíprocamente, si $q=2^k$, entonces $p<2^k$ se escribe con $k$ cifras binarias, $p=(c_{k-1}\dots c_0)_2$, y $p/2^k=(0{,}c_{k-1}\dots c_0)_2$.
\end{proof}

Por ejemplo, $0{,}1=1/10$ tiene denominador $10=2\cdot5$, de modo que su desarrollo binario es infinito: $0{,}1=(0{,}0\overline{0011})_2$. Ningún ordenador binario almacena $0{,}1$ exactamente.

\section{Binario, hexadecimal y bytes}

Un \emph{bit} es una cifra binaria y un \emph{byte} son 8 bits, que representan $2^8=256$ valores. Como $16=2^4$, cada cifra \emph{hexadecimal} ($0$--$9$, $\texttt{A}$--$\texttt{F}$) corresponde a un grupo de 4 bits, y un byte se escribe con dos cifras hexadecimales, de $\texttt{0x00}$ a $\texttt{0xFF}$. Para convertir entre binario y hexadecimal basta agrupar los bits de cuatro en cuatro desde la derecha: $(1011\,0110)_2=\texttt{0xB6}=182$.

Para cantidades de memoria se usan dos sistemas de prefijos: los decimales del SI ($1\ \text{kB}=10^3$ bytes, $1\ \text{GB}=10^9$ bytes) y los binarios de la IEC ($1\ \text{KiB}=2^{10}$ bytes, $1\ \text{GiB}=2^{30}$ bytes). La diferencia es de un $7{,}4\,\%$ en el caso de los gigas.

\section{Enteros con signo: complemento a dos}

\begin{definicion}[Complemento a dos]\label{def:c2}
Una palabra de $n$ bits $w=(a_{n-1}\dots a_0)$ representa, \emph{sin signo}, el entero $U(w)=\sum_{k=0}^{n-1}a_k2^k$ y, en \emph{complemento a dos}, el entero
\[
T(w)=-a_{n-1}2^{n-1}+\sum_{k=0}^{n-2}a_k2^k=U(w)-a_{n-1}2^n.
\]
\end{definicion}

\begin{teorema}\label{teo:c2}
Para $n\ge2$:
\begin{enumerate}
  \item $T$ es una biyección entre las palabras de $n$ bits y los enteros del intervalo $[-2^{n-1},\,2^{n-1}-1]$; $T(w)<0$ si y solo si $a_{n-1}=1$;
  \item $T(w)$ y $U(w)$ difieren en un múltiplo de $2^n$;
  \item si $\bar w$ es la palabra con todos los bits invertidos, $T(\bar w)+1=-T(w)$; en particular, el opuesto de $T(w)$ se obtiene invirtiendo los bits y sumando~1, y el resultado es representable en $n$ bits salvo si $T(w)=-2^{n-1}$.
\end{enumerate}
\end{teorema}

\begin{proof}
(i) Por la \cref{prop:cifras}, $U$ es una biyección sobre $\{0,\dots,2^n-1\}$. Si $a_{n-1}=0$, $U(w)<2^{n-1}$ y $T(w)=U(w)\in[0,2^{n-1}-1]$. Si $a_{n-1}=1$, $U(w)\in[2^{n-1},2^n-1]$ y $T(w)=U(w)-2^n\in[-2^{n-1},-1]$. Cada mitad se aplica biyectivamente sobre su imagen, y las imágenes son disjuntas y cubren el intervalo. (ii) Es la igualdad $T(w)=U(w)-a_{n-1}2^n$. (iii) $U(w)+U(\bar w)=\sum_k2^k=2^n-1$, luego $T(w)+T(\bar w)=2^n-1-(a_{n-1}+\bar a_{n-1})2^n=-1$, porque exactamente uno de los dos bits de signo vale~1. Así, $T(\bar w)+1=-T(w)$ como igualdad de enteros; si $T(w)=-2^{n-1}$, el valor $-T(w)=2^{n-1}$ no está en el rango de~(i).
\end{proof}

\begin{ejemplo}\label{ej:c2}
Con $n=8$ el rango es $[-128,127]$. Como $5=(0000\,0101)_2$, se tiene $-5=(1111\,1010)_2+1=(1111\,1011)_2$; en efecto, $-128+64+32+16+8+2+1=-5$.
\end{ejemplo}

Por el \cref{teo:c2}(ii), la suma de enteros con signo se realiza con el mismo circuito que la suma sin signo, descartando el acarreo final: el resultado es correcto salvo que salga del rango, situación que se llama \emph{desbordamiento} (\cref{eje:6-desbordamiento}).

\section{Coma flotante}

\begin{definicion}[Número en coma flotante normalizado]\label{def:flotante}
Un formato binario de coma flotante con $e$ bits de exponente, sesgo $\beta=2^{e-1}-1$ y $p$ bits de fracción representa la palabra $(s,E,f)$, con $1\le E\le2^e-2$, como
\[
x=(-1)^s\,(1{,}f_1f_2\dots f_p)_2\cdot2^{E-\beta}.
\]
Los valores extremos de $E$ se reservan para el cero, los números subnormales, los infinitos y \texttt{NaN}~\cite{ieee754}.
\end{definicion}

\begin{table}[h]
\centering\small
\caption{Formatos binarios de coma flotante habituales. El épsilon es $\varepsilon=2^{-p}$, la distancia de~1 al siguiente número representable.}\label{tab:formatos}
\begin{tabular}{@{}lcccccc@{}}
\toprule
Formato & Bits & Signo & Exponente $e$ & Fracción $p$ & $\varepsilon=2^{-p}$ & Máximo finito \\
\midrule
fp64 (doble) & 64 & 1 & 11 & 52 & $2{,}2\cdot10^{-16}$ & $1{,}8\cdot10^{308}$ \\
fp32 (simple) & 32 & 1 & 8 & 23 & $1{,}2\cdot10^{-7}$ & $3{,}4\cdot10^{38}$ \\
fp16 (media) & 16 & 1 & 5 & 10 & $9{,}8\cdot10^{-4}$ & $65\,504$ \\
bf16 (bfloat16) & 16 & 1 & 8 & 7 & $7{,}8\cdot10^{-3}$ & $3{,}4\cdot10^{38}$ \\
\bottomrule
\end{tabular}
\end{table}

El formato bfloat16 conserva los 8 bits de exponente de fp32 y recorta la fracción a 7 bits: tiene el mismo rango que fp32 con mucha menos precisión, y sus 16 bits son exactamente los 16 bits superiores de un fp32.

\begin{proposicion}[Error de redondeo]\label{prop:redondeo}
Sea $\fl(x)$ el número del formato más cercano a $x$. Si $x$ está en el rango de los números normalizados, entonces
\[
\fl(x)=x(1+\delta)\qquad\text{con}\qquad\abs{\delta}\le u=2^{-(p+1)}=\frac\varepsilon2.
\]
\end{proposicion}

La demostración es el \cref{eje:6-redondeo}. El número $u$ se llama \emph{unidad de redondeo}: el error \emph{relativo} es pequeño, pero el absoluto crece con $\abs{x}$, porque la distancia entre números consecutivos en $[2^k,2^{k+1})$ es $2^{k-p}$~\cite{goldberg}.

\begin{observacion}[Errores frecuentes]
\begin{enumerate}[label=(\alph*)]
\item Confundir GB y GiB, o kB y KiB.
\item Olvidar que el rango del complemento a dos es asimétrico: en 8 bits, $-(-128)$ desborda.
\item Comparar números en coma flotante con igualdad exacta: $0{,}1+0{,}2\neq0{,}3$ en doble precisión, porque ninguno de los tres es representable.
\item Creer que fp16 y bf16 son intercambiables: el primero tiene más precisión y el segundo mucho más rango.
\end{enumerate}
\end{observacion}

\begin{observacion}[Conexión con el curso]
Los tokenizadores BPE a nivel de byte (módulo~12) parten de un vocabulario inicial de exactamente $2^8=256$ símbolos, uno por byte, y el texto se convierte antes en bytes mediante UTF-8 (\cref{eje:6-utf8}). El entrenamiento en precisión mixta con bf16 o fp16 (módulo~11) y la cuantización a enteros de 8 o 4 bits (módulo~17) son elecciones entre rango, precisión y memoria que se analizan con las herramientas de este tema.
\end{observacion}

\begin{lecturas}
\item Bryant y O'Hallaron~\cite{csapp}, §2.1 «Information Storage» (hexadecimal, bytes), §2.2 «Integer Representations» (complemento a dos, extensión de signo), §2.3 «Integer Arithmetic» (desbordamiento) y §2.4 «Floating Point».
\item Goldberg~\cite{goldberg}, sección sobre errores de redondeo (épsilon de la máquina y unidad de redondeo).
\item IEEE 754~\cite{ieee754}, para la definición de los formatos y del redondeo al par más cercano.
\end{lecturas}

\section*{Ejercicios del tema 6}
\addcontentsline{toc}{section}{Ejercicios del tema 6}

\begin{ejercicio}{1}{eje:6-a-decimal}
Pase a base~10: $(1011\,0110)_2$, $\texttt{0x2F}$, $\texttt{0xFF}$ y $(777)_8$.
\fuente{adaptado de Bryant y O'Hallaron~\cite{csapp}, §2.1, problemas de práctica de conversión.}
\end{ejercicio}

\begin{ejercicio}{1}{eje:6-contar}
(a) ¿Cuántos valores distintos pueden representarse con 8, 16 y 32 bits? (b) ¿Cuál es el rango de un entero sin signo de 8 bits y el de uno con signo en complemento a dos? (c) ¿Cuántas cifras hexadecimales ocupa una palabra de 32 bits?
\end{ejercicio}

\begin{ejercicio}{1}{eje:6-vf}
Razone si son verdaderas o falsas: (a) $0{,}1$ es exactamente representable en fp32; (b) bf16 y fp32 tienen aproximadamente el mismo máximo finito; (c) en complemento a dos, el bit más significativo indica el signo; (d) $0{,}625$ es exactamente representable en fp16.
\end{ejercicio}

\begin{ejercicio}{2}{eje:6-a-binario}
Escriba en binario y en hexadecimal: $200$, $1000$ y $4095$. Compruebe cada resultado con el esquema de Horner.
\fuente{adaptado de Bryant y O'Hallaron~\cite{csapp}, §2.1.}
\end{ejercicio}

\begin{ejercicio}{2}{eje:6-fracciones}
Escriba en base~2: (a) $0{,}625$; (b) $13/16$; (c) $1/3$; (d) $0{,}1$. Explique, con la \cref{prop:fraccion-binaria}, cuáles tienen desarrollo finito.
\end{ejercicio}

\begin{ejercicio}{2}{eje:6-c2}
Con palabras de 8 bits en complemento a dos: (a) represente $-1$, $-100$ y $127$; (b) interprete $(1000\,0001)_2$ y $(1111\,0110)_2$; (c) calcule $100+50$ con aritmética de 8 bits e interprete el resultado.
\fuente{adaptado de Bryant y O'Hallaron~\cite{csapp}, §2.2--2.3.}
\end{ejercicio}

\begin{ejercicio}{2}{eje:6-utf8}
UTF-8 codifica los puntos de código de \texttt{U+0080} a \texttt{U+07FF} con dos bytes, y los de \texttt{U+0800} a \texttt{U+FFFF} con tres bytes, según las plantillas
\[
\texttt{110xxxxx 10xxxxxx}\qquad\text{y}\qquad\texttt{1110xxxx 10xxxxxx 10xxxxxx},
\]
donde las $x$ son los bits del punto de código, completados con ceros a la izquierda. Codifique «é» (\texttt{U+00E9}) y «€» (\texttt{U+20AC}) y dé los bytes en hexadecimal.
\end{ejercicio}

\begin{ejercicio}{2}{eje:6-digitos}
Usando la \cref{prop:cifras}: (a) ¿cuántas cifras binarias tiene $10^9$? (b) ¿cuántas cifras decimales tiene $2^{100}$? (Use $\log_{10}2\approx0{,}30103$.) (c) ¿Cuántas cifras hexadecimales tiene $2^{32}-1$?
\end{ejercicio}

\begin{ejercicio}{2}{eje:6-decodificar}
(a) Interprete las palabras fp32 $\texttt{0x3F800000}$ y $\texttt{0xC0A00000}$. (b) Codifique $0{,}75$ y $3$ en bf16 y dé el resultado en hexadecimal.
\fuente{adaptado de Bryant y O'Hallaron~\cite{csapp}, §2.4.}
\end{ejercicio}

\begin{ejercicio}{2}{eje:6-espaciado}
En bf16 ($p=7$) con redondeo al más cercano (en caso de empate, al de última cifra par): (a) ¿cuál es la distancia entre números consecutivos en $[2^k,2^{k+1})$? (b) Calcule $\fl(256+1)$, $\fl(\fl(256+1)+1)$ y $\fl(256+\fl(1+1))$. (c) ¿Cuál es el menor entero positivo no representable en bf16? ¿Y en fp32? (d) Compruebe que el máximo finito de fp16 es $65\,504$.
\end{ejercicio}

\begin{ejercicio}{2}{eje:6-memoria}
Un modelo tiene $7\cdot10^{9}$ parámetros. Calcule la memoria que ocupan sus pesos en fp32, bf16, int8 y con 4 bits por parámetro, en GB y, para fp32, también en GiB.
\end{ejercicio}

\begin{ejercicio}{3}{eje:6-divisibilidad}
Demuestre que un entero es divisible por~3 (respectivamente, por~9) si y solo si la suma de sus cifras decimales lo es. Enuncie y demuestre el criterio análogo para la divisibilidad por~15 en base~16. (Indicación: \cref{eje:3-diferencia-potencias}.)
\end{ejercicio}

\begin{ejercicio}{3}{eje:6-extension}
Demuestre que la \emph{extensión de signo} conserva el valor en complemento a dos: si $w$ es una palabra de $n$ bits y $w'$ la de $n+1$ bits que se obtiene repitiendo a la izquierda el bit de signo, entonces $T(w')=T(w)$.
\fuente{adaptado de Bryant y O'Hallaron~\cite{csapp}, §2.2.}
\end{ejercicio}

\begin{ejercicio}{3}{eje:6-decimal-finito}
Demuestre que una fracción irreducible $p/q$ tiene desarrollo decimal finito si y solo si $q=2^a5^b$ con $a,b\in\N_0$.
\end{ejercicio}

\begin{ejercicio}{3}{eje:6-desbordamiento}
Sean $x,y$ enteros en $[-2^{n-1},2^{n-1}-1]$ y sea $s$ el resultado de la suma en complemento a dos, es decir, el único entero del rango que difiere de $x+y$ en un múltiplo de $2^n$. Demuestre: (a) si $x$ e $y$ tienen signos distintos, no hay desbordamiento ($s=x+y$); (b) hay desbordamiento si y solo si $x$ e $y$ tienen el mismo signo y $s$ tiene el signo contrario.
\fuente{adaptado de Bryant y O'Hallaron~\cite{csapp}, §2.3.}
\end{ejercicio}

\begin{ejercicio}{3}{eje:6-redondeo}
Demuestre la \cref{prop:redondeo}: si $2^k\le\abs{x}<2^{k+1}$ y el formato tiene $p$ bits de fracción, entonces $\abs{\fl(x)-x}\le2^{k-p-1}$ y $\abs{\fl(x)-x}/\abs{x}\le2^{-(p+1)}$.
\fuente{adaptado de Goldberg~\cite{goldberg}, sección sobre errores de redondeo.}
\end{ejercicio}
